\section{Gaps and Opportunities}
\label{sec:gaps-opportunities}

\subsection{Architecture Gap: Experimentability}
\label{subsec:gap-experimentability}
The prior Reading the Reader theses each delivered a working prototype, but each was tightly coupled and was extended by copying the previous frontend rather than by plugging into a shared platform \cite{refsgaard2024rtr, pereira2024typography, palinec2024reactive}.
This coupling hinders the systematic comparison of interventions and decision strategies that the project ultimately requires.
At the same time, empirical findings show that intervention design materially affects recovery and user preference \cite{jensen2025context}, which makes such comparison scientifically important rather than merely convenient.

These coupling problems point to a clear opportunity.
Designing the adaptation engine around hot-swappable intervention modules and auditable logging would enable systematic A/B testing, replay-based evaluation, real-time researcher override, and the transparent experimental comparison of interventions that the prior prototypes could not support.

\subsection{Sensing Gap: Research-Grade versus Accessible Hardware}
\label{subsec:gap-sensing}
The market survey in Sec.~\ref{sec:market-landscape} exposes a split in the sensing layer.
Research-grade eye trackers deliver the sampling rate and precision needed for reliable fixation and saccade extraction, but they are closed and costly \cite{tobii2025fusion, srresearch2025eyelink, gazepoint2025gp3, pupillabs2025neon}, whereas the more accessible reading and accessibility tools do not expose a gaze signal that can drive adaptation \cite{orcam2025read3, microsoft2025immersive, amazon2025kindle}.
A platform tied to one tracker therefore inherits both its cost and its constraints.

The opportunity is to place sensing behind a stable adapter interface so that the eye tracker is one interchangeable implementation among several.
This decouples the reading runtime from any single vendor, allows research-grade hardware to be used where available while leaving room for lower-cost or webcam-based sensing, and keeps the door open to the multimodal signals (including pupil size) envisaged for the project \cite{dtucompute2025rtr}.

\subsection{Integration Gap: External and AI-Driven Decision Modules}
\label{subsec:gap-integration}
None of the surveyed products expose a pluggable decision layer; each bundles fixed logic for when and how to act on the reader.
This is at odds with the project's vision, in which decisions about adaptation may eventually be made by external, AI-driven providers rather than hard-coded rules \cite{dtucompute2025rtr}.
Without a defined integration boundary, every new decision strategy would require modifying the core application.

The opportunity is to define a decision-strategy contract that the core application depends on, so that built-in rules and external or AI-based providers are interchangeable behind the same interface.
This makes AI support an architectural property of the platform rather than a feature implemented end-to-end in this thesis, and it lets future contributors supply decision providers without touching the sensing, analysis, or intervention layers.
